% GENERATED FILE. Edit canonical lessons, then run npm run generate:books.
% canonical-source-hash: fnv1a64:38910e2ffa1dc177
% canonical-lessons: TE-C180-pradhana, TE-C180-mottam, TE-C180-prati, TE-C180-gulabi-rangu, TE-C180-narinja-rangu

\chapter{Main, Total, Every, Pink, Orange}
\label{ch:te-main-total-every-pink-orange}
\hlchaptermodality{180}

\begin{chapteropening}
\textbf{By the end of this chapter:} \emph{I can say main, total, every, pink and orange.}

Complete the last lesson of chapter 180: I can say main, total, every, pink and orange.
\end{chapteropening}

\section[\texorpdfstring{\te{ప్రధాన} (pradhāna)}{ప్రధాన (pradhāna)}]{\te{ప్రధాన} (pradhāna) — main, chief}

\label{lesson:TE-C180-pradhana}



\begin{quote}
Before the new one: say the Telugu for expensive, then the Telugu for cheap.
\end{quote}

\subsection*{You'll want to know: \te{ప్రధాన}}
\textbf{\te{ప్రధాన}} — \emph{pradhāna} — \textquotedblleft{}main, chief\textquotedblright{}.

\begin{grammarlens}[title={What you've built}]
One more word for this chapter.
\end{grammarlens}

\subsection*{Your turn}
\begin{itemize}
\raggedright
  \item \emph{Say it:} \emph{pradhāna}
  \item \emph{Say it:} \emph{pradhāna}, once more
  \item \emph{Say it:} say \emph{pradhāna}
  \item \emph{Recall:} say the Telugu for expensive, then the Telugu for cheap, then say \emph{pradhāna} again
\end{itemize}

\begin{tcolorbox}[breakable,colback=teal!4,colframe=teal!35!black,title={Before you move on}]
What does \te{ప్రధాన} mean? (\textquotedblleft{}Main, chief\textquotedblright{}.) Say it once more.
\end{tcolorbox}
\section[\texorpdfstring{\te{మొత్తం} (mottaṁ)}{మొత్తం (mottaṁ)}]{\te{మొత్తం} (mottaṁ) — total, whole; the total}

\label{lesson:TE-C180-mottam}



\begin{quote}
Before the new one: say the Telugu for cheap, then the Telugu for main.
\end{quote}

\subsection*{You'll want to know: \te{మొత్తం}}
\textbf{\te{మొత్తం}} — \emph{mottaṁ} — \textquotedblleft{}total, whole; the total\textquotedblright{}.

\begin{grammarlens}[title={What you've built}]
One more word for this chapter.
\end{grammarlens}

\subsection*{Your turn}
\begin{itemize}
\raggedright
  \item \emph{Say it:} \emph{mottaṁ}
  \item \emph{Say it:} \emph{mottaṁ}, once more
  \item \emph{Say it:} say \emph{mottaṁ}
  \item \emph{Recall:} say the Telugu for cheap, then the Telugu for main, then say \emph{mottaṁ} again
\end{itemize}

\begin{tcolorbox}[breakable,colback=teal!4,colframe=teal!35!black,title={Before you move on}]
What does \te{మొత్తం} mean? (\textquotedblleft{}Total, whole; the total\textquotedblright{}.) Say it once more.
\end{tcolorbox}
\section[\texorpdfstring{\te{ప్రతి} (prati)}{ప్రతి (prati)}]{\te{ప్రతి} (prati) — every, each}

\label{lesson:TE-C180-prati}



\begin{quote}
Before the new one: say the Telugu for main, then the Telugu for total.
\end{quote}

\subsection*{You'll want to know: \te{ప్రతి}}
\textbf{\te{ప్రతి}} — \emph{prati} — \textquotedblleft{}every, each\textquotedblright{}.

\begin{grammarlens}[title={What you've built}]
One more word for this chapter.
\end{grammarlens}

\subsection*{Your turn}
\begin{itemize}
\raggedright
  \item \emph{Say it:} \emph{prati}
  \item \emph{Say it:} \emph{prati}, once more
  \item \emph{Say it:} say \emph{prati}
  \item \emph{Recall:} say the Telugu for main, then the Telugu for total, then say \emph{prati} again
\end{itemize}

\begin{tcolorbox}[breakable,colback=teal!4,colframe=teal!35!black,title={Before you move on}]
What does \te{ప్రతి} mean? (\textquotedblleft{}Every, each\textquotedblright{}.) Say it once more.
\end{tcolorbox}
\section[\texorpdfstring{\te{గులాబీ} \te{రంగు} (gulābī raṅgu)}{గులాబీ రంగు (gulābī raṅgu)}]{\te{గులాబీ} \te{రంగు} (gulābī raṅgu) — pink}

\label{lesson:TE-C180-gulabi-rangu}



\begin{quote}
Before the new one: say the Telugu for total, then the Telugu for every.
\end{quote}

\subsection*{You'll want to know: \te{గులాబీ} \te{రంగు}}
\textbf{\te{గులాబీ} \te{రంగు}} — \emph{gulābī raṅgu} — \textquotedblleft{}pink\textquotedblright{}.

\begin{grammarlens}[title={What you've built}]
One more word for this chapter.
\end{grammarlens}

\subsection*{Your turn}
\begin{itemize}
\raggedright
  \item \emph{Say it:} \emph{gulābī raṅgu}
  \item \emph{Say it:} \emph{gulābī raṅgu}, once more
  \item \emph{Say it:} say \emph{gulābī raṅgu}
  \item \emph{Recall:} say the Telugu for total, then the Telugu for every, then say \emph{gulābī raṅgu} again
\end{itemize}

\begin{tcolorbox}[breakable,colback=teal!4,colframe=teal!35!black,title={Before you move on}]
What does \te{గులాబీ} \te{రంగు} mean? (\textquotedblleft{}Pink\textquotedblright{}.) Say it once more.
\end{tcolorbox}
\section[\texorpdfstring{\te{నారింజ} \te{రంగు} (nāriñja raṅgu)}{నారింజ రంగు (nāriñja raṅgu)}]{\te{నారింజ} \te{రంగు} (nāriñja raṅgu) — orange (colour)}

\label{lesson:TE-C180-narinja-rangu}



\begin{quote}
Before the new one: say the Telugu for every, then the Telugu for pink.
\end{quote}

\subsection*{You'll want to know: \te{నారింజ} \te{రంగు}}
\textbf{\te{నారింజ} \te{రంగు}} — \emph{nāriñja raṅgu} — \textquotedblleft{}orange (colour)\textquotedblright{}.

\begin{grammarlens}[title={What you've built}]
One more word for this chapter.
\end{grammarlens}

\subsection*{Your turn}
\begin{itemize}
\raggedright
  \item \emph{Say it:} \emph{nāriñja raṅgu}
  \item \emph{Say it:} \emph{nāriñja raṅgu}, once more
  \item \emph{Say it:} say \emph{nāriñja raṅgu}
  \item \emph{Recall:} say the Telugu for every, then the Telugu for pink, then say \emph{nāriñja raṅgu} again
\end{itemize}

\begin{tcolorbox}[breakable,colback=teal!4,colframe=teal!35!black,title={Before you move on}]
What does \te{నారింజ} \te{రంగు} mean? (\textquotedblleft{}Orange (colour)\textquotedblright{}.) Say it once more.
\end{tcolorbox}
